\begin{solapartat}
\figura[0.45]{sol_fig1.png}

\punts{0,45} Com que el punt on hem de calcular el camp està situat entre les dues càrregues i aquestes tenen signe diferent, el camp total serà la suma dels valors absoluts dels camps de cadascuna de les càrregues, la direcció serà la de la recta que uneix les dues càrregues i el sentit serà el que va de la càrrega positiva cap a la negativa.

\punts{0,25} $\left|\vec{E}_1\right| = \left|k\,\dfrac{q_1}{d_1^2}\right| = 8,99\times 10^{9}\,\dfrac{5\times 10^{-6}}{(0,03)^2} = 4,99\times 10^{7}\ \mathrm{N/C}$

% DUBTE: la pauta escriu q_1/d_1^2 a l'expressió de E_2 (hauria de ser q_2/d_2^2); els valors numèrics sí que són els de la càrrega 2.
\punts{0,25} $\left|\vec{E}_2\right| = \left|k\,\dfrac{q_1}{d_1^2}\right| = 8,99\times 10^{9}\,\dfrac{7\times 10^{-6}}{(0,07)^2} = 1,28\times 10^{7}\ \mathrm{N/C}$

I com els dos camps són paral·lels, el mòdul de la suma és la suma de mòduls:

\punts{0,3} $\left|\vec{E}_{Total}\right| = \left|\vec{E}_1\right| + \left|\vec{E}_2\right| = 6,28\times 10^{7}\ \mathrm{N/C}$
\end{solapartat}

\begin{solapartat}
\punts{0,25} $V_P = V_1 + V_2 = 0$

\punts{0,4} $0 = -8,99\times 10^{9}\,\dfrac{5\times 10^{-6}}{\mathrm{d}} + 8,99\times 10^{9}\,\dfrac{7\times 10^{-6}}{0,1-\mathrm{d}}$

\punts{0,4} $\dfrac{5}{\mathrm{d}} = \dfrac{7}{0,1-\mathrm{d}} \Longrightarrow \dfrac{\mathrm{d}}{5} = \dfrac{0,1-\mathrm{d}}{7} \Longrightarrow 12d = 0,5 \Longrightarrow d = 0,0417\ \mathrm{m}$ \quad \punts{0,2}
\end{solapartat}
